\begin{abstract}
Medical image segmentation in clinical settings is hindered by stringent
data-privacy regulations that prevent centralising patient images for model
training.
Federated Learning (FL) offers a compelling remedy: clients train on local
data and share only model updates.
However, real-world medical data exhibit severe label heterogeneity---both
across institutions (non-IID distributions) and within institutions
(annotation scarcity)---which substantially degrades the performance of
vanilla FL methods.

We present a comprehensive FL framework for 2-D medical image segmentation
that (i) evaluates four aggregation strategies---FedAvg, FedProx, FedNova,
and FedPer---under controlled Dirichlet non-IID partitioning
($\alpha \in \{0.1, 0.5, 1.0\}$); and (ii) introduces \textbf{FedAL}, a
novel federated active-learning protocol that couples entropy-based query
selection with communication-efficient aggregation to minimise labelling
cost across distributed clients.
All methods share a common UNet backbone
(channels $[32,64,128,256,512]$) with a compound Dice$+$CE loss, and are
evaluated on four public benchmarks: ACDC (cardiac MRI, 3 classes),
BUSI (breast ultrasound, 1 class), FUGC (gastrointestinal polyp, 2 classes),
and TN3K (thyroid nodule, 1 class).

Experimental results demonstrate that FedNova consistently outperforms
FedAvg on heterogeneous partitions, reducing the DSC gap by up to
\textbf{X.X\%} on ACDC ($\alpha=0.1$).
FedPer achieves the best DSC on BUSI, highlighting the benefit of
personalized decoders for ultrasound modalities.
FedAL reaches competitive DSC with $\leq 20\%$ of the labelling budget
required by fully supervised FL.
Our system ranked \textbf{1st} on the FUGC 2025 challenge leaderboard,
validating the effectiveness of the proposed framework.
\end{abstract}
